% Fixed130 complete independent insertion.
% Analytic 2-integral splitting of ad(T) on the genuine tame S3
% sector; strict elementary chain reductions and Hecke top obstruction.

\subsection{Tame-idempotent splitting and an analytic integral
adjoint Fox--Smith certificate in the $S_3$ sector}
\label{subsec:dyadic-s3-adjoint-tame-splitting}

The coefficient-specific matrix computation in
Theorem~\ref{thm:dyadic-s3-marked-adjoint-jacobian}
admits a stronger structural proof.  For the fixed
integral $S_3$ representation $T=\mathbf Z_2^2$,
the adjoint module $\operatorname{End}(T)$
has a canonical order-three tame idempotent.
Its invariant and noninvariant summands are
both explicitly rank two, and the latter
is \emph{integrally isomorphic} to the natural
coefficient representation treated by
Theorem~\ref{thm:dyadic-s3-natural-strict-retract}.
The invariant summand is an induced
unramified-quadratic permutation lattice.
This splits the full $8\times16$ marked
Jacobian \emph{before} elementary reduction,
giving a symbolic proof of its unique
factor $2$ and an actual generator
of the top obstruction.

Write, as in \eqref{eq:dyadic-s3-integral-matrices},
\[
S=\begin{pmatrix}0&1\\1&0\end{pmatrix},
\qquad
U=\begin{pmatrix}-1&-1\\1&0\end{pmatrix},
\quad
M=\operatorname{End}_{\mathbf Z_2}(T).
\]
Let $\operatorname{Ad}_g(A)=gAg^{-1}$.
The order-three averaging operator is
\begin{equation}\label{eq:dyadic-s3-adjoint-tame-projector}
\boxed{\quad
\Pi_0=\tfrac13
\bigl(I+\operatorname{Ad}_U+
\operatorname{Ad}_{U^2}\bigr),\quad
M_0=\operatorname{im}\Pi_0,\quad
M_1=\ker\Pi_0.
\quad}
\end{equation}
The idempotent commutes with $\operatorname{Ad}_S$,
because $SUS^{-1}=U^2$; it also commutes
with the whole fixed $S_3$ action since
$\langle U\rangle\triangleleft S_3$.

\begin{theorem}[A unimodular tame splitting of $\operatorname{ad}T$]
\label{thm:dyadic-s3-adjoint-unimodular-splitting}
Define two $2$-adic integral matrices
\begin{equation}\label{eq:dyadic-s3-adjoint-vw-matrices}
 V=(I-\Pi_0)E_{12}
 =\frac13\begin{pmatrix}-1&1\\2&1\end{pmatrix},
 \qquad
 W=SVS^{-1}
 =\frac13\begin{pmatrix}1&2\\1&-1\end{pmatrix}.
\end{equation}
Then
\begin{equation}\label{eq:dyadic-s3-adjoint-basis}
\boxed{\quad
M_0=\mathbf Z_2 U\oplus\mathbf Z_2 U^2,
\qquad
M_1=\mathbf Z_2 V\oplus\mathbf Z_2 W,
\qquad
M=M_0\oplus M_1.
\quad}
\end{equation}
In the row-major basis
$(E_{11},E_{12},E_{21},E_{22})$,
the change-of-basis matrix with
columns $(U,U^2,V,W)$ is
\begin{equation}\label{eq:dyadic-s3-adjoint-unimodular-basis}
\mathsf B=
\begin{pmatrix}
-1&0&-1/3&1/3\\
-1&1&1/3&2/3\\
1&-1&2/3&1/3\\
0&-1&1/3&-1/3
\end{pmatrix},
\qquad
\boxed{\det\mathsf B=1\in\mathbf Z_2^\times}.
\end{equation}
The $S_3$ action is, in these
two ordered rank-two blocks,
\begin{equation}\label{eq:dyadic-s3-adjoint-s3-block-actions}
\begin{array}{c|cc}
& M_0:(U,U^2)&M_1:(V,W)\\ \hline
\operatorname{Ad}_U
&I_2&
\begin{pmatrix}-1&-1\\1&0\end{pmatrix}\\[6pt]
\operatorname{Ad}_S
&\begin{pmatrix}0&1\\1&0\end{pmatrix}&
\begin{pmatrix}0&1\\1&0\end{pmatrix}.
\end{array}
\end{equation}
Consequently, as \emph{actual
integral} $\mathbf Z_2[S_3]$-modules,
\begin{equation}\label{eq:dyadic-s3-adjoint-induced-plus-natural}
\boxed{\quad
\operatorname{ad}T
\cong M_0\oplus M_1
\cong
\operatorname{Ind}_{G_{K_2}}^\Gamma\mathbf Z_2
\oplus T,
\qquad K_2=\mathbf Q_2(\zeta_3).
\quad}
\end{equation}
Here the induction is from the
index-two normal subgroup $G_{K_2}$
and uses its two cosets as a
\emph{permutation lattice};
it is not induction from the
index-three subgroup $H$.
\end{theorem}

\begin{proof}
The operator $\Pi_0$ is the
standard Reynolds idempotent
for the normal order-three
subgroup; $3$ is invertible in
$\mathbf Z_2$.  Its invariant
part is the centralizer
$\mathbf Z_2[U]$, which is free
of rank two with basis $(U,U^2)$
because $U^2+U+I=0$.
Evaluating $(I-\Pi_0)E_{12}$
gives the displayed $V$ and $W$.
The row-major matrix of these
four vectors is
\eqref{eq:dyadic-s3-adjoint-unimodular-basis}.
Multiplying its last two
columns by $3$ gives an
integer matrix of determinant $9$;
hence $\det\mathsf B=9/9=1$.
Thus $(U,U^2,V,W)$ is a
$\mathbf Z_2$ basis, and the
projector proves the direct sum
\eqref{eq:dyadic-s3-adjoint-basis}.

Conjugation by $S$ swaps
$U$ and $U^2$, and swaps $V,W$
by definition.  Conjugation by
$U$ fixes its centralizer
and, by direct multiplication,
\[
 UVU^{-1}=-V+W,\qquad
 UWU^{-1}=-V.
\]
These are exactly the two
matrices in
\eqref{eq:dyadic-s3-adjoint-s3-block-actions}.
The $M_1$ matrices are precisely
the natural representation
\eqref{eq:dyadic-s3-integral-matrices};
the $M_0$ matrices describe the
permutation action on the two
cosets of the normal subgroup
$\langle U\rangle$.
The preimage of $\langle U\rangle$
in $\Gamma$ is $G_{K_2}$,
so this permutation lattice is
$\operatorname{Ind}_{G_{K_2}}^\Gamma
\mathbf Z_2$.
This proves the integral
module isomorphism.
\end{proof}

\begin{theorem}[An analytic proof of the fixed128 adjoint Smith factors]
\label{thm:dyadic-s3-adjoint-analytic-smith}
In the coefficient basis
\eqref{eq:dyadic-s3-adjoint-basis}
the complete full-marked
three-term adjoint \emph{coefficient}
complex of fixed128 is the direct
sum of a $M_0$ complex and the
norm-zero natural complex on $M_1$.
The latter is the strict
$M_1[-1]$ retract of
Theorem~\ref{thm:dyadic-s3-natural-strict-retract}
via the canonical identification
$M_1\cong T$.

On the tame-invariant rank-two block
$M_0$ choose $S_M=\begin{psmallmatrix}
0&1\\1&0\end{psmallmatrix}$ in
the $(U,U^2)$ basis.
Since $\tau$ acts trivially
and $x_0,x_1$ act trivially,
the full marked differentials are
\begin{equation}\label{eq:dyadic-s3-tame-invariant-differentials}
\boxed{
\begin{aligned}
d_0^0(v)
&=((S_M-I)v,\ 0,\ 0,\ 0),\\
d_0^1(a,b,c,d)
&=\bigl((S_M-2I)b,\
  3b+2c+(S_M-I)d\bigr).
\end{aligned}}
\end{equation}
The operator $L=S_M-2I$ has
determinant $3$ and is a
$\mathbf Z_2$-unit.
The strict invertible target
change
\begin{equation}\label{eq:dyadic-s3-tame-invariant-row-change}
(u,w)\longmapsto
(u,w-3L^{-1}u)
\end{equation}
changes the degree-one
differential to
\begin{equation}\label{eq:dyadic-s3-tame-invariant-split-d1}
d_0^1(a,b,c,d)
\longmapsto
\bigl(Lb,\ 2c+(S_M-I)d\bigr).
\end{equation}
It therefore splits off the
contractible tame block
$M_0\xrightarrow{L}M_0$
in degrees $1$ and $2$,
leaving the three-generator
integral Fox complex
\begin{equation}\label{eq:dyadic-s3-tame-invariant-fox-core}
\boxed{\quad
M_0\xrightarrow{(S_M-I,0,0)^t}M_0^3
\xrightarrow{(a,c,d)\mapsto
2c+(S_M-I)d}M_0.
\quad}
\end{equation}

Put
\[
e=(1,0)^t,\qquad
e_+=(1,1)^t,\qquad
f=(-1,1)^t=(S_M-I)e.
\]
Both $(e_+,e)$ and $(f,e)$
are unimodular bases of $M_0$.
The following elementary
unit operations are sufficient
to reduce
\eqref{eq:dyadic-s3-tame-invariant-fox-core}:
\begin{enumerate}[label=\textnormal{(\roman*)}]
\item $e\mapsto f$ in $d^0$
splits off a degree-zero--one
unit pair, leaving $e_+$
as a generator of $H^0$;
\item the input $d=e$ maps
to $f$ in degree two,
giving one degree-one--two
unit pair;
\item $c=e$ maps to $2e$
modulo the image of $f$,
giving one surviving
two-primary elementary
pair $\mathbf Z_2\xrightarrow{2}\mathbf Z_2$;
\item the remaining three
degree-one coordinates
are freely generated
by a $\sigma$-value $e$,
by $(c,d)=(f,-2e)$,
and by $(c,d)=(0,e_+)$.
\end{enumerate}
Thus the tame-invariant block has
the strict finite homotopy type
\begin{equation}\label{eq:dyadic-s3-tame-invariant-elementary}
\boxed{\quad
\mathcal J^\bullet(M_0)
\simeq_{\mathbf Z_2}
\mathbf Z_2[0]\oplus
\mathbf Z_2^3[-1]\oplus
[\mathbf Z_2\xrightarrow{2}
 \mathbf Z_2]_{[1,2]}.
\quad}
\end{equation}
Together with
$\mathcal J^\bullet(M_1)\simeq M_1[-1]$,
the full adjoint marked complex has
the \emph{analytically derived}
integral elementary form
\begin{equation}\label{eq:dyadic-s3-adjoint-elementary-from-tame}
\boxed{\quad
\mathcal J^\bullet_\Gamma(M)
\simeq_{\mathbf Z_2}
\mathbf Z_2[0]\oplus
\mathbf Z_2^5[-1]\oplus
[\mathbf Z_2\xrightarrow{2}
 \mathbf Z_2]_{[1,2]}.
\quad}
\end{equation}
In particular its degree-one
Jacobian has \emph{seven unit
Smith factors and exactly one
factor of $2$}, recovering
\eqref{eq:dyadic-s3-jacobian-smith}
without extracting a numerical
mod-$4$ minor.  Its degree-zero
Jacobian has three unit factors,
and its cohomology is precisely
\[
(H^0,H^1,H^2)(\Gamma,M)=
(\mathbf Z_2,\mathbf Z_2^5,\mathbf F_2).
\]
\end{theorem}

\begin{proof}
All marked differential matrices
are linear in the coefficient
representation, so the
$\mathbf Z_2[S_3]$ splitting
of Theorem~\ref{thm:dyadic-s3-adjoint-unimodular-splitting}
gives a strict splitting of
their cochain complex.
On $M_1$, the exact identity
$I+\operatorname{Ad}_U+
\operatorname{Ad}_{U^2}=0$
gives the natural norm-zero
proof of fixed129, with its
explicit strong contraction.

On $M_0$, the affine coefficient
action of $\tau$ is trivial.
For arbitrary crossed-derivation
values $(a,b,c,d)$,
the tame relation gives
$(S_M-2I)b$.  In the wild
relation the pro-$2$ projections
of $x_i\tau$ act as the identity
on each translation, so
$u_i=x_i\tau$,
$d_0=d_g=\tau$, and
$c_0=h_c=1$.  Since
$g_0=\sigma^2$ has trivial
linear part on $M_0$,
conjugation of a translation
by $g_0$ is trivial.
Thus
$h_0=2c+4b$,
while $u_1^{-1}x_1^\sigma$
adds $-d-b+S_M^{-1}d$.
As $S_M^{-1}=S_M$,
the wild differential is
$3b+2c+(S_M-I)d$,
proving
\eqref{eq:dyadic-s3-tame-invariant-differentials}.
The row change
\eqref{eq:dyadic-s3-tame-invariant-row-change}
then cancels the tame unit.
The remaining three-generator
complex is literally the
pro-$2$ marked Fox cochain
complex of
Theorem~\ref{thm:dyadic-strict-nielsen-fox-chain}
in these coefficients.

For the elementary reduction,
$(S_M-I)e=f$ and
$(S_M-I)e_+=0$.
The image of $d^0$ is the
primitive rank-one line
$\mathbf Z_2 f$ in the
$\sigma$ coordinate, so
one integral unit pair splits.
In degree one the $d=e$
coordinate maps to the
primitive target vector $f$.
Eliminating that row changes
the image of $c=e$ to
$2e$ in the complementary
free target quotient.
The other input coordinates
give exactly the three free
cocycles listed in (iv),
as can be checked by
$2c+(S_M-I)d=0$.
All changes use unimodular
$\mathbf Z_2$ bases; this
proves
\eqref{eq:dyadic-s3-tame-invariant-elementary}.
The direct sum with $M_1[-1]$
then gives
\eqref{eq:dyadic-s3-adjoint-elementary-from-tame}.
Counting its cancelled
degree-one--two blocks gives
two tame units plus one
Fox unit from $M_0$, and
four unit degree-one--two
blocks from $M_1$, hence
seven units and the
remaining factor $2$.
The degree-zero--one units
are one from $M_0$ and
two from $M_1$, hence three.
\end{proof}

\begin{corollary}[Complete marked adjoint cocycle and obstruction coordinates]
\label{cor:dyadic-s3-adjoint-explicit-h1-h2}
In the marked $S_3$ coefficient basis
$(U,U^2,V,W)$, every class in
$H^1(\Gamma,\operatorname{ad}T)$ has
a unique representative modulo
coboundaries of the form
\begin{equation}\label{eq:dyadic-s3-adjoint-normalized-cocycles}
\boxed{
\begin{aligned}
c(\sigma)&=\lambda e\in M_0,\\
c(\tau)&=0,\\
c(x_0)&=\mu f+t,\\
c(x_1)&=-2\mu e+\nu e_++2S_M t,
\end{aligned}
\qquad
\lambda,\mu,\nu\in\mathbf Z_2,\quad
t\in M_1\cong\mathbf Z_2^2.
}
\end{equation}
Here $e,e_+,f$ are the vectors
in the $(U,U^2)$ basis from
Theorem~\ref{thm:dyadic-s3-adjoint-analytic-smith},
while $S_M$ acts by the same swap
matrix on either block.
The parameters
$(\lambda,\mu,\nu,t)$ form
an exact integral
$\mathbf Z_2^5$ basis
of the global adjoint
$H^1$.

The top adjoint cohomology
has exactly one nonzero
$\mathbf F_2$ class:
in the reduced $M_0$
Fox target it is represented
by $e$ modulo
$2M_0+(S_M-I)M_0$.
The ordinary matrix trace
maps this class to
$\operatorname{tr}(U)=-1$
modulo two, hence
gives its canonical
nonzero scalar detector.
The trace annihilates
$M_1$ and therefore
locates the top
obstruction \emph{entirely}
in the tame-invariant
summand $M_0$.
\end{corollary}

\begin{proof}
The unit tame row forces
$c(\tau)=0$ for an actual
cocycle in either summand.
On $M_1$, the strong
normalization of
Theorem~\ref{thm:dyadic-s3-natural-strict-retract}
gives $(c(\sigma),c(x_0),c(x_1))
=(0,t,2S_Mt)$.
On $M_0$, the reduced Fox
equation $2c+(S_M-I)d=0$
forces $c=\mu f$ and
$d=-2\mu e+\nu e_+$.
The $\sigma$ coordinate is
defined modulo the primitive
coboundary line $\mathbf Z_2 f$,
and so has the unique
representative $\lambda e$.
These descriptions produce
the displayed five parameters.
The top Fox quotient is
\[
M_0/(2M_0+(S_M-I)M_0)
\cong\mathbf F_2,
\]
generated by $e$.
On $M_0=\mathbf Z_2[U]$,
matrix trace sends $U$ and
$U^2$ to $-1$, hence
annihilates $(S_M-I)M_0$,
and also annihilates
$2M_0$ modulo two.
Its value on $e=U$
is nonzero.  Since the
matrices $V,W$ are traceless,
the detector vanishes on
$M_1$ and proves the
final assertion.
\end{proof}

\begin{theorem}[Integral Shapiro decomposition and top-degree Hecke selection]
\label{thm:dyadic-s3-adjoint-shapiro-hecke}
Put $K_2=\mathbf Q_2(\zeta_3)$,
$F=\mathbf Q_2(\sqrt[3]2)$,
$L=FK_2$, and let
$N=G_L\subset H=G_F\subset\Gamma$.
For the fixed $S_3$ action on
$M=\operatorname{ad}T$,
the module decomposition of
Theorem~\ref{thm:dyadic-s3-adjoint-unimodular-splitting}
gives the following
\emph{integral} Shapiro and
restriction identifications:
\begin{equation}\label{eq:dyadic-s3-adjoint-shapiro-global-h}
\boxed{
\begin{aligned}
R\Gamma(\Gamma,M_0)&\simeq
R\Gamma(G_{K_2},\mathbf Z_2),\\
R\Gamma(\Gamma,M_1)&\simeq
R\Gamma(\Gamma,T),\\
R\Gamma(H,M_0)&\simeq
R\Gamma(N,\mathbf Z_2),\\
R\Gamma(H,M_1)&\simeq
R\Gamma(N,\mathbf Z_2).
\end{aligned}}
\end{equation}
The cohomology ledgers are
\begin{equation}\label{eq:dyadic-s3-adjoint-shapiro-ledgers}
\begin{array}{c|ccc}
& H^0&H^1&H^2\\ \hline
\Gamma,M_0&\mathbf Z_2&\mathbf Z_2^3&\mathbf F_2\\
\Gamma,M_1&0&\mathbf Z_2^2&0\\
H,M_0&\mathbf Z_2&\mathbf Z_2^7&\mathbf F_2\\
H,M_1&\mathbf Z_2&\mathbf Z_2^7&\mathbf F_2 .
\end{array}
\end{equation}
The index-three Sylow idempotent
$e_H=\tfrac13\operatorname{res}\operatorname{cor}$
preserves $M_0\oplus M_1$ and
acts on $H^2(H,M)$ as
\begin{equation}\label{eq:dyadic-s3-adjoint-top-hecke-projection}
\boxed{\quad
H^2(H,M)=\mathbf F_2\oplus\mathbf F_2,
\qquad
e_H^{(2)}=
\begin{pmatrix}1&0\\0&0\end{pmatrix}
\quad}
\end{equation}
with the first coordinate
the tame-invariant component.
Likewise its degree-one
image/kernel decomposition is
\[
\begin{array}{c|cc}
&\operatorname{rank}\operatorname{im}e_H&
\operatorname{rank}\ker e_H\\ \hline
H^1(H,M_0)&3&4\\
H^1(H,M_1)&2&5.
\end{array}
\]
Thus the one surviving global
two-primary obstruction in
fixed128 is the
\emph{odd-index restriction}
of the trivial-coefficient
obstruction over
$K_2=\mathbf Q_2(\zeta_3)$,
while the entire top class
from the natural $M_1$
summand is killed by
the Hecke/Sylow projector.
\end{theorem}

\begin{proof}
The first two Shapiro identities
follow directly from
\eqref{eq:dyadic-s3-adjoint-induced-plus-natural}.
On restriction to $H$,
the generator $\sigma$ acts
on both $M_0$ and $M_1$
by the swap matrix, while
the wild generators act
trivially.  The index-two
kernel of this action
inside $H$ is $N=G_L$.
Thus each restricted module
is the permutation-induced
lattice $\operatorname{Ind}_N^H
\mathbf Z_2$, proving the
remaining two Shapiro identities.
The degree-one
and degree-two cohomology
of a trivial integral
$\mathbf Z_2$ lattice
over a dyadic local
absolute Galois group
is given by local class
field theory and Tate
duality:
\[
H^0(G_D,\mathbf Z_2)=\mathbf Z_2,
\quad H^1(G_D,\mathbf Z_2)
=\mathbf Z_2^{[D:\mathbf Q_2]+1},
\quad H^2(G_D,\mathbf Z_2)=\mathbf F_2
\]
for $D=K_2,L$, since
both contain exactly $\mu_2$
among the $2$-power roots
of unity.  Combined with
the natural-$T$ calculation
of fixed128, these give
\eqref{eq:dyadic-s3-adjoint-shapiro-ledgers}.

For $M_0$ the ordinary Mackey
identification
\[
\operatorname{Res}_H^\Gamma
\operatorname{Ind}_{G_{K_2}}^\Gamma
\mathbf Z_2
\simeq
\operatorname{Ind}_N^H\mathbf Z_2
\]
holds because $H G_{K_2}=\Gamma$
and $H\cap G_{K_2}=N$.
Under Shapiro, restriction
$\Gamma\to H$ becomes
the usual restriction
$G_{K_2}\to N=G_L$.
Since $[L:K_2]=3$ and
$3$ is a unit modulo $2$,
this restriction is an
isomorphism on the
one-dimensional
$H^2(-,\mathbf Z_2)$:
its local Tate-dual
effect is multiplication
by the odd degree three
on the $\mu_2$ invariant.
Thus the idempotent
acts as identity on
the first $H^2$ component.
For $M_1$, global
$H^2(\Gamma,M_1)=0$ by
the norm-zero calculation,
so the same idempotent
annihilates the second
$H^2$ component.
This gives
\eqref{eq:dyadic-s3-adjoint-top-hecke-projection}.
The degree-one ranks follow
from the Shapiro ledger
and the usual
$\operatorname{cor}\operatorname{res}=3$
splitting, separately on
each summand.
\end{proof}

\begin{warning}[Scope of the analytic tame decomposition]
\label{warn:dyadic-s3-adjoint-tame-splitting-scope}
The tame projector $\Pi_0$
splits the \emph{fixed}
$S_3$ adjoint coefficient
lattice and any family
whose entire coefficient
action preserves the
normal $C_3$-isotypic
decomposition.  On an
arbitrary universal framed
deformation of
$\bar\rho:G_{\mathbf Q_2}\to
\mathrm{GL}_2(\mathbf F_2)$,
the wild generator matrices
need not commute with $\Pi_0$.
Therefore the decomposition
must \emph{not} be asserted
as an equivariant direct
sum on the full deformation
ring without additional
tame-isotypic invariance
hypotheses.

The analytic computation
upgrades fixed128 from
a numerical Smith
certificate to explicit
integral cochain
elementary divisors
and a top-degree Hecke
selection \emph{for this
specific adjoint lattice}.
It does not select
canonical minimal generators
of the rank-five
Demu\v{s}kin group
$G_F(2)$, or give
the universal Schreier
relation-syzygy chain
isomorphism required
for arbitrary non-pro-$2$
residual families.
\end{warning}
